\section{Related Work}
\label{sec:related}

\paragraph{Composed image retrieval.}
TIRG established image-text composition as a learned retrieval problem and introduced a gated residual fusion operator~\cite{vo2019tirg}. FashionIQ later provided natural-language feedback for fashion search~\cite{wufashioniq}, while CIRR extended evaluation to open-domain images~\cite{liucirr}. PinPoint broadens evaluation with explicit negatives, paraphrases, and multi-image queries~\cite{mahadev2026pinpoint}. Early supervised methods learned the composed query through joint visual-semantic matching~\cite{chen2020jvsm}, visiolinguistic attention~\cite{chen2020val}, modality-agnostic fusion~\cite{dodds2020maaf}, or fashion-specific compositional transitions~\cite{yu2020curlingnet}. Other approaches used residual composition~\cite{shin2021rtic}, dual objectives~\cite{kim2021dcnet}, comprehensive linguistic-visual fusion~\cite{wen2021clvc}, or attribute-aware contrastive learning~\cite{tian2022aacl}. Later systems introduced content-style modulation~\cite{lee2021cosmo}, semantic attention~\cite{jandial2022sac}, explicit/implicit target matching~\cite{delmas2022artemis}, and multi-grained uncertainty regularization~\cite{chen2022mur}. CLIP4CIR strengthened supervised composition with task-oriented CLIP features and contrastive learning~\cite{baldrati2022clip4cir}. DQU-CIR moved fusion to the input level through textual and visual query unification~\cite{wen2024dqu}. ConText-CIR aligns textual concepts with relevant image content~\cite{xing2025context}, CCIN neutralizes compositional conflicts~\cite{tian2025ccin}, and MAPNet models multi-schema interactions and noisy negatives~\cite{shi2025mapnet}. CoLLM uses large language models to create and encode multimodal training queries~\cite{huynh2025collm}; ReCALL recalibrates retrieval ability after MLLM adaptation~\cite{yang2026recall}; MCoT-MVS selects multi-level evidence through multimodal reasoning~\cite{ge2026mcot}. FashionMV introduces product-level multi-view retrieval and ProCIR, providing the FashionGen protocol used in this work~\cite{yuan2026fashionmv}.

\paragraph{Data-efficient and zero-shot CIR.}
Because annotated CIR triplets are costly, a large body of work removes or reduces task-specific supervision. Pic2Word maps an image into the language token space~\cite{saito2023pic2word}, and SEARLE learns the corresponding pseudo word through textual inversion~\cite{baldrati2023searle}. LinCIR trains a projection using language-only self-supervision~\cite{gu2024lincir}. KEDs enriches pseudo-word representations with an external database and textual concept alignment~\cite{suo2024keds}, while CompoDiff edits image features with a latent diffusion model~\cite{gu2024compodiff}. RTD adapts the text encoder after pretraining~\cite{byun2025rtd}, while HIT learns hierarchical pseudo words with text-conditioned filtering~\cite{li2025hit}. SLERP and text-anchored tuning combine image and text representations geometrically~\cite{jang2024slerp}. Visual Delta Generator creates semi-supervised triplets with a multimodal generator~\cite{jang2024vdg}. Recent approaches also predict missing target content~\cite{tang2025worldmodel}, reason before retrieval~\cite{tang2025reason}, generate imagined proxies~\cite{li2025imagine}, augment retrieval with composed image generation~\cite{wang2025generative}, or learn under noisy triplets~\cite{li2026conesep}. These methods improve the contents or training of one query representation. \method instead composes already trained retrieval endpoints and is therefore complementary to advances in endpoint design.

\paragraph{Ensembling and representation compatibility.}
Score averaging is a strong retrieval baseline because independently trained models often make complementary errors. Its usual form is diagonal: each query encoder is paired only with its own gallery encoder. \method asks a different question. If two systems inhabit aligned coordinates, can the query function learned by one system retrieve against the gallery geometry learned by another? The resulting cross paths preserve the standard nearest-neighbor interface, but expose directional interactions hidden by diagonal averaging. Our compute-matched diagonal ensemble separates this effect from the benefit of running two endpoints, and the coordinate-scrambling experiment separates genuine compatibility from coincidental score aggregation.
